% Formalización pública del resultado autoral de polarización causal.

\section{Campo de polarización causal y pérdida de libertad radial}
\label{sec:md-polarizacion-causal-definicion}

Sea \((\mathcal M,g)\) una geometría esféricamente simétrica escrita en
coordenadas de Painlevé--Gullstrand,
\begin{equation}
 \mathrm ds^2
 =-c^2\,\mathrm dt^2
  +\bigl(\mathrm dr+v(r)\,\mathrm dt\bigr)^2
  +r^2\,\mathrm d\Omega_2^2,
 \label{eq:md-polarizacion-metrica-pg}
\end{equation}
donde \(v(r)\geq0\) mide la inclinación radial de los conos causales respecto
de la congruencia temporal seleccionada. El campo adimensional
\begin{equation}
 \mathcal P(r):=\frac{v(r)}{c}
 \label{eq:md-polarizacion-parametro}
\end{equation}
es la intensidad de polarización causal. Sobre la distribución espacial
ortogonal a la congruencia temporal se define el tensor
\begin{equation}
 \Pi(r)=\mathcal P(r)^2\,\widehat r\otimes\widehat r,
 \label{eq:md-polarizacion-tensor}
\end{equation}
cuya dirección principal es radial y cuyas direcciones tangenciales
permanecen degeneradas. Esta definición no introduce una dirección absoluta:
la dirección principal está determinada covariantemente por el gradiente del
radio areal.

Las generatrices radiales nulas de \eqref{eq:md-polarizacion-metrica-pg}
satisfacen
\begin{equation}
 \frac{\mathrm dr}{\mathrm dt}=-v(r)\pm c
 =c\bigl(-\mathcal P(r)\pm1\bigr).
 \label{eq:md-polarizacion-generatrices}
\end{equation}
Por tanto, el conjunto de signos radialmente accesibles cambia de cardinalidad
al atravesar \(\mathcal P=1\).

\begin{theorem}[Dirección causal obligada de caída]
\label{thm:md-polarizacion-caida-obligada}
En la geometría \eqref{eq:md-polarizacion-metrica-pg} se verifican las
siguientes tres alternativas exhaustivas:
\begin{enumerate}
 \item si \(0\leq\mathcal P<1\), una generatriz radial futura aumenta \(r\)
 y la otra lo disminuye;
 \item si \(\mathcal P=1\), la generatriz saliente satisface
 \(\mathrm dr/\mathrm dt=0\) y define el horizonte;
 \item si \(\mathcal P>1\), ambas generatrices radiales futuras satisfacen
 \(\mathrm dr/\mathrm dt<0\).
\end{enumerate}
En el tercer régimen el grado radial no desaparece de la variedad, pero deja
de constituir una elección espacial bidireccional: todo vector causal futuro
posee componente dirigida hacia radios areales menores. El interior es, en
este sentido exacto, una región de caída causal obligada.
\end{theorem}

\begin{proof}
Las tres conclusiones se obtienen comparando los dos valores
\(-\mathcal P\pm1\) de \eqref{eq:md-polarizacion-generatrices} con cero. Para
\(\mathcal P>1\), incluso la rama con signo positivo es negativa. La
conclusión concierne al cono causal y no a una elección particular de
coordenada temporal.
\end{proof}

\section{Realización de Schwarzschild y gradiente colectivo de caída}
\label{sec:md-polarizacion-schwarzschild}

Para una masa total \(M\), la realización de Schwarzschild en la carta
\eqref{eq:md-polarizacion-metrica-pg} viene dada por
\begin{equation}
 v(r)=c\sqrt{\frac{R}{r}},
 \qquad
 R=\frac{2GM}{c^2},
 \qquad
 \mathcal P(r)^2=\frac{R}{r}.
 \label{eq:md-polarizacion-schwarzschild}
\end{equation}
El horizonte queda caracterizado simultáneamente por
\begin{equation}
 r=R,
 \qquad
 \mathcal P=1,
 \qquad
 g^{\mu\nu}\partial_\mu r\,\partial_\nu r=0.
 \label{eq:md-polarizacion-horizonte}
\end{equation}
La aceleración radial de la congruencia de caída se recupera como gradiente
de la intensidad cinética del campo:
\begin{equation}
 a_r=-\frac12\frac{\mathrm d}{\mathrm dr}v(r)^2
 =-\frac{GM}{r^2}.
 \label{eq:md-polarizacion-gradiente}
\end{equation}
Así, la gravitación radial se expresa como gradiente colectivo de la caída y
no como una fuerza añadida después de fijar las trayectorias.

\begin{corollary}[Macropolarización causal del interior gravitatorio]
\label{cor:md-macropolarizacion-interior}
En la realización \eqref{eq:md-polarizacion-schwarzschild}, la región
\(r<R\) satisface \(\mathcal P>1\). El agujero negro constituye una
macropolarización causal: la dirección principal de \(\Pi\) rebasa el umbral
unitario y alinea ambas familias radiales futuras hacia la disminución del
radio areal.
\end{corollary}

\section{Correspondencia con los \texorpdfstring{\(3\)}{3} regímenes del TRIT}
\label{sec:md-polarizacion-trit}

El lector geométrico del TRIT se realiza mediante
\begin{equation}
 \tau(r)=\operatorname{sgn}\!\bigl(1-\mathcal P(r)^2\bigr)
 \in\{+1,0,-1\}.
 \label{eq:md-polarizacion-lector-trit}
\end{equation}
La correspondencia es
\begin{equation}
 \begin{array}{c|c|c|c}
 \tau & \mathcal P & \text{régimen geométrico} &
 \text{orientación radial futura}\\ \hline
 +1 & \mathcal P<1 & \text{exterior} & \text{entrada y salida}\\
 0  & \mathcal P=1 & \text{horizonte} & \text{generatriz saliente nula}\\
 -1 & \mathcal P>1 & \text{interior} & \text{caída obligada}
 \end{array}
 \label{eq:md-polarizacion-tabla-trit}
\end{equation}
La involución temporal del TRIT intercambia los regímenes exterior e interior
y fija el horizonte. El cambio \(+1\leftrightarrow-1\) no elimina la memoria
de la ruta: modifica su orientación causal y conserva la frontera
\(\tau=0\) como superficie de empalme.

\Needspace{30\baselineskip}
\section{Cierre local de materia y alineamiento gravitatorio colectivo}
\label{sec:md-polarizacion-materia}

Sea \(\gamma_a\) una sección material persistente con proyector radial
\(\widehat r_a\otimes\widehat r_a\) y peso de cierre \(w_a\geq0\). Para un
conjunto finito de secciones se define el parámetro de orden
\begin{equation}
 \Pi_N
 =\frac{1}{\sum_{a=1}^{N}w_a}
  \sum_{a=1}^{N}w_a\,
  \widehat r_a\otimes\widehat r_a.
 \label{eq:md-polarizacion-colectiva}
\end{equation}
Una partícula individual realiza un cierre local persistente de ruta. El
interior gravitatorio aparece cuando el campo colectivo fija una dirección
principal causal para todas las rutas futuras. Por ello no se identifican dos
objetos matemáticamente distintos:
\begin{equation}
\begin{aligned}
 \text{partícula}
 &=\text{sección material persistente},\\
 \text{agujero negro}
 &=\text{alineamiento causal macroscópico por encima del umbral}.
\end{aligned}
 \label{eq:md-polarizacion-distincion}
\end{equation}

La secuencia causal completa de esta residencia queda fijada por
\begin{equation}
\begin{aligned}
 \text{estado APP--TRIT--TPK}
 &\longrightarrow
 \text{sección material}
 \longrightarrow
 \text{campo de polarización }\Pi\\
 &\longrightarrow
 \mathcal P=1
 \longrightarrow
 \text{horizonte}
 \longrightarrow
 \mathcal P>1
 \longrightarrow
 \text{caída causal obligada}.
\end{aligned}
 \label{eq:md-polarizacion-cadena}
\end{equation}
El capítulo siguiente prolonga esta estructura hacia la inversión de hoja,
la purificación avanzada--retardada, la geometría de Kruskal y la información
de horizonte.

\section{Condiciones geométricas de la macropolarización causal: radio areal, inclinación de conos, horizonte y caída interior}
\label{sec:md-polarizacion-control}

El resultado exige simultáneamente que el radio utilizado sea areal, que
\(v(r)\) incline los conos según \eqref{eq:md-polarizacion-metrica-pg}, que
el horizonte satisfaga \(\mathcal P=1\) y que las dos pendientes de
\eqref{eq:md-polarizacion-generatrices} sean negativas en el interior. Estas
identidades distinguen la pérdida efectiva de libertad radial de una mera
metáfora dimensional y proporcionan un control directo de la
macropolarización causal.
